% Modelo criado por Chaiane Bitelo para o Simpósio de Pós-Graduação em Design do PPDESDI - UERJ
% jun/2020 - Versão 1


\documentclass[a4paper,11pt,oneside,onecolumn,final]{article}


\input{estrutura} % Chama o arquivo estrutura.tex onde estão os pacotes e principais ajustes de formatação 

\begin{document}
\setlength{\parindent}{0pt} % configura todo o documento para não ter identação


% =======================================================
% PRIMEIRA PÁGINA
% =======================================================
\begin{titlepage}

% ----- CABEÇALHO DA PRIMEIRA PÁGINA ----------------------- 
\textsc{spgd2020} \hfill \break
\textsc{6º simpósio de pós-graduação em design da esdi} \hfill \break
 Rio de Janeiro, 04, 05, 06, 07 de novembro de 2020
%-----------------------------------------------------------

\vspace{10mm} % insere espaço entre o parágrafo anterior e o próximo  

% ------- TÍTULO DO ARTIGO -----------------   
\Huge O título para artigo completo do SPGD 2020 deve ser composto em até três linhas: referência para formatação
% ------------------------------------------

\vspace{5mm} % insere espaço entre o parágrafo anterior e o próximo

% ------ 
% Título em Inglês
% -------

\huge \textit{Title for SPGD 2019 paper in English}

\vspace{10mm} % espaço entre os parágrafos

% ------ 
% Autores
% -------
\normalsize
SOBRENOME, Nome; SOBRENOME, Nome; SOBRENOME, Nome;SOBRENOME, Nome;

\vspace{5mm} % espaço entre os parágrafos

% ------ 
% Resumo e palavras-chave em Português
% -------


\textbf{RESUMO:} Apresentamos aqui as diretrizes para formatação e submissão de artigos completos para o SPGD 2019. O resumo deverá conter de 100 a 150 palavras, em parágrafo único, o suficiente para expor o tema geral, a abordagem da pesquisa e seus principais resultados. Este documento é referência para formatação do artigo, que será publicado conforme enviado pelos autores. A primeira página deverá conter o cabeçalho, o título em português e sua versão em inglês, o resumo e o abstract, com palavras-chave e keywords. Repita as palavras-chave utilizadas no resumo expandido e utilize os respectivos termos em inglês aqui sugeridos.

\vspace{5mm} % espaço entre os parágrafos

\textbf{Palavras-chave}: Arte e visualidade; Criatividade e processos de criação; Cultura e sociedade; Design da informação; Design de interação; Design de moda; Design de negócios; Design de produto; Design e antropologia; Design e arquitetura; Design e políticas públicas; Design e serviços; Design e sustentabilidade; Design e tecnologia; Design gráfico; Educação e pedagogia do design; Ergonomia e usabilidade; Gestão do design; História do design; Imagem e comunicação; Pesquisa e metodologia do design; Teoria do design.

\vspace{10mm} % espaço entre os parágrafos

% ------ 
% Resumo e palavras-chave em Inglês
% -------


\textbf{ABSTRACT:} \textit{These are the guidelines for formatting and submission of papers for SPGD 2019. The abstract should contain 100-150 words, enough to expose the general theme, the research approach and its main results. This document is a reference to article formatting, which will be reproduced as submitted by the authors. The first page should contain the header, the Portuguese title and its English version, authors' names, the abstract in Portuguese and in English with keywords. Repeat the keywords used in your expanded abstract and the terms in English suggested below}

\vspace{5mm} %5mm vertical space

\textbf{Keywords}: \textit{Art and visuality; Business design; Creativity and creation processes; Culture and society; Design and anthropology; Design and architecture; Design and public policies; Design and services; Design and sustainability; Design and technology; Design education and pedagogy; Design history; Design management; Design theory; Ergonomics and usability; Fashion design; Graphic design; Image and communication; Information design; Interaction design; Product design; Research and design methodology.}

\end{titlepage}

% ============================================================= 
% CORPO DO ARTIGO - a partir da segunda página
% =============================================================

\setlength{\parskip}{5pt} % configura espacejamento entre parágrafos
\renewcommand{\baselinestretch}{1.15} % configura espaçamento entre linhas


\section{— Orientações gerais}

Os artigos devem apresentar pesquisas de mestrandos, doutorandos, e egressos do PPDESDI. Estudantes de graduação da ESDI também estão convidados a participar. Os resumos expandidos submetidos e aceitos deverão ser apresentados oralmente durante o Simpósio. Além disso, deverão ser apresentados na forma de vídeo com duração máxima de cinco minutos e elaborados conforme template disponível em \url{www.even3.com.br/spgd_2020}. O artigo completo será enviado após o simpósio. Para que o artigo seja publicado nos Anais será necessário que pelo menos um dos autores esteja inscrito no evento. Havendo dúvida, entre em contato através do e-mail: esdi.spgd@gmail.com.

\section{— Formatação do artigo}

Artigos submetidos para o 5º Simpósio de Pós-Graduação em Design da ESDI – SPGD 2020 devem ter entre \textbf{6 e 12 páginas}, não numeradas, incluindo resumo, notas de rodapé, legendas, figuras, tabelas e referências. Os arquivos deverão ser enviados em formato PDF. O tamanho máximo do arquivo não deve ultrapassar 10 megabytes. O  arquivo deve ser identificado como ''SPGD2020 nome do seu artigo''. 

\subsection{Estrutura}

A primeira página do artigo deve conter: título em português, título em inglês, resumo em português, palavras-chave em português, resumo em inglês e palavras-chave em inglês. No título, use maiúsculas apenas na primeira palavra, em nomes próprios e siglas. Repita as cinco palavras-chave utilizadas no seu resumo expandido. Informações complementares sobre os autores deverão ser inseridas ao final do artigo em ''Notas sobre os autores''.

\subsection{Expressões lógicas}

Para inclusão de expressões lógicas ou matemáticas, utilize o pacote {amsmath,amssymb} 

\subsection{Cabeçalhos}

O cabeçalho da segunda página em diante deve conter unicamente os sobrenomes dos autores separados por ponto e vírgula e o título do artigo em português. Essa informação deve ser incluída no arquivo chamado '' estrutura.tex ''. O cabeçalho deverá ser composto de em apenas uma linha. Assim, títulos e subtítulos longos deverão ter as últimas palavras suprimidas. As supressões devem ser substituídas por reticências. 

\subsection{Citações}

Para citar um autor no corpo do seu texto use “aspas”, sem itálico. Citações que excedam três linhas devem se destacar em parágrafo próprio, como a seguir:

\begin{quote}
    Citações com mais de três linhas devem ser formatadas em bloco utilizando o comando ''quote'', destacadas do texto. Não se esqueça de incluir a indicação da página de onde o trecho foi retirado, indicando no colchete junto ao comando ''cite''. Não há necessidade de usar aspas ou itálico, uma vez que o destaque está dado pelo recuo e espaços \cite[p. 1]{bonsiepe_teoria_1993}.

\end{quote}
Citações no corpo do texto seguem a norma ABNT NBR 10520:2002. pelo sistema “autor-data”. 

Exemplos: 

(1)
Segundo \citeonline[p. 9]{bonsiepe_metodologia_1984}, ''gerar alternativas técnicas básicas não é sinônimo de \textit{brainstorming} e do registro da corrente de todas as possibilidades.

(2)
Nesse sentido, é possível afirmar que ''gerar alternativas técnicas básicas não é sinônimo de \textit{brainstorming} e do registro da corrente de todas as possibilidades \cite[p. 9]{bonsiepe_metodologia_1984}.  


\subsection{Figuras e Tabelas}

Caso utilize tabelas e figuras, desenhos, fotografias e diagramas, utilize os modelos abaixo.

\textbf{Figuras}: nas figuras, caso haja crédito, insira o nome do autor ou detentor dos direitos da imagem, identificando-o na legenda, assim como a fonte de onde foi retirada a imagem, como mostrado na Figura \ref{fig:PredioEsdi} a seguir:

\begin{figure}[ht]
    \includegraphics[width=12cm]{PredioEsdiPB.png} % o tamanho pode ser alterado substituindo o valor em centímetros conforme necessidade. Para a imagem ocupar a mesma largura o texto use o comando \texwidth
    \caption{Título da Fotografia ou da Imagem. Crédito: Fotógrafo ou dono da foto (se houver)}
    \label{fig:PredioEsdi}
\end{figure}

Para que sua imagem ocupe a largura da página, use o comando conforme o modelo abaixo na Figura \ref{fig:PredioEsdi2}.

\begin{figure}[ht]
    \includegraphics[width=\textwidth]{PredioEsdiPB.png} 
    \caption{Título da Fotografia ou da Imagem. Crédito: Fotógrafo ou dono da foto (se houver)}
    \label{fig:PredioEsdi2} % crie aqui o label dando um nome curto para sua figura e use o comando \ref no corpo do texto para referenciar
\end{figure}

Para citar as figuras, use como este exemplo onde a Figura \ref{fig:PredioEsdi} está sendo referência com o uso dos comandos label e ref.


\textbf{Tabelas}: no caso das tabelas, utilize o modelo abaixo.

\begin{table}[!ht]
  \caption{Exemplo de tabela.}
    \begin{tabular}{|+l|^r|^r|^r|^r|}
    \hline
    \rowstyle{\bfseries}
        Coluna A & Coluna B & Coluna C & Coluna D \\
    \hline
          Texto 1 & B1 & C1 & D1\\
    \hline
          Texto 2 & B2 & C2 & D2\\        
    \hline
          Texto 3 & B3 & C3 & D3\\
    \hline
          Texto 4 & B4 & C4 & D4\\      
    \hline
  \end{tabular}\hfill
    \caption*{Fonte: Sobrenome, ano, p. x. (se houver).}
    \label{mais.rotulo} % crie aqui o label dando um nome curto para sua tabela e use o comando \ref no corpo do texto para referenciar
\end{table}

% -----------------------
% CONCLUSÕES DO ARTIGO
% ----------------------


\section{— Conclusões}

\lipsum[1-3]

% -----------------------
% Agradecimentos
% ----------------------

\section*{Agradecimentos}

Não se esqueça de mencionar os órgãos de fomento (Faperj, Capes, CNPq ou outros), o PPDESDI, os professores, a UERJ, as instituições com que têm vínculos e que permitem a realização do trabalho, os participantes da pesquisa.

% ---------------
% Referências 
% --------------

\bibliography{bibliografia}


% ----------------------------------
% NOTAS SOBRE OS AUTORES
% ----------------------------------

.  .  .  .  .  .  .  .  .  .  .  .  .  .  .  .  . 

\textbf{Notas sobre os autores}

Sobrenome, Nome; Titulação; Título provisório da Dissertação/Tese; Orientador; Ano previsto para defesa; link para Currículo Lattes 
\url{http://lattes.cnpq.br}
\url{nome.sobrenome@provedor.com}

\vspace{5mm} % espaço de 5mm entre os parágrafos 

Sobrenome, Nome; Titulação; Título provisório da Dissertação/Tese; Orientador; Ano previsto para defesa; link para Currículo Lattes 
\url{http://lattes.cnpq.br}
\url{nome.sobrenome@provedor.com}

\vspace{5mm} % espaço de 5mm entre os parágrafos 

Sobrenome, Nome; Titulação; Título provisório da Dissertação/Tese; Orientador; Ano previsto para defesa; link para Currículo Lattes 
\url{http://lattes.cnpq.br}
\url{nome.sobrenome@provedor.com}

\end{document}
